\IfFileExists{stacks-project.cls}{%
\documentclass{stacks-project}
}{%
\documentclass{amsart}
}

% For dealing with references we use the comment environment
\usepackage{verbatim}
\newenvironment{reference}{\comment}{\endcomment}
%\newenvironment{reference}{}{}
\newenvironment{slogan}{\comment}{\endcomment}
\newenvironment{history}{\comment}{\endcomment}

% For commutative diagrams we use Xy-pic
\usepackage[all]{xy}

% We use 2cell for 2-commutative diagrams.
\xyoption{2cell}
\UseAllTwocells

% We use multicol for the list of chapters between chapters
\usepackage{multicol}

% This is generally recommended for better output
\usepackage{lmodern}
\usepackage[T1]{fontenc}

% For cross-file-references
\usepackage{xr-hyper}

% Package for hypertext links:
\usepackage{hyperref}

% For any local file, say "hello.tex" you want to link to please
% use \externaldocument[hello-]{hello}
\externaldocument[introduction-]{introduction}
\externaldocument[conventions-]{conventions}
\externaldocument[sets-]{sets}
\externaldocument[categories-]{categories}
\externaldocument[topology-]{topology}
\externaldocument[sheaves-]{sheaves}
\externaldocument[sites-]{sites}
\externaldocument[stacks-]{stacks}
\externaldocument[fields-]{fields}
\externaldocument[algebra-]{algebra}
\externaldocument[brauer-]{brauer}
\externaldocument[homology-]{homology}
\externaldocument[derived-]{derived}
\externaldocument[simplicial-]{simplicial}
\externaldocument[more-algebra-]{more-algebra}
\externaldocument[smoothing-]{smoothing}
\externaldocument[modules-]{modules}
\externaldocument[sites-modules-]{sites-modules}
\externaldocument[injectives-]{injectives}
\externaldocument[cohomology-]{cohomology}
\externaldocument[sites-cohomology-]{sites-cohomology}
\externaldocument[dga-]{dga}
\externaldocument[dpa-]{dpa}
\externaldocument[sdga-]{sdga}
\externaldocument[hypercovering-]{hypercovering}
\externaldocument[schemes-]{schemes}
\externaldocument[constructions-]{constructions}
\externaldocument[properties-]{properties}
\externaldocument[morphisms-]{morphisms}
\externaldocument[coherent-]{coherent}
\externaldocument[divisors-]{divisors}
\externaldocument[limits-]{limits}
\externaldocument[varieties-]{varieties}
\externaldocument[topologies-]{topologies}
\externaldocument[descent-]{descent}
\externaldocument[perfect-]{perfect}
\externaldocument[more-morphisms-]{more-morphisms}
\externaldocument[flat-]{flat}
\externaldocument[groupoids-]{groupoids}
\externaldocument[more-groupoids-]{more-groupoids}
\externaldocument[etale-]{etale}
\externaldocument[chow-]{chow}
\externaldocument[intersection-]{intersection}
\externaldocument[pic-]{pic}
\externaldocument[weil-]{weil}
\externaldocument[adequate-]{adequate}
\externaldocument[dualizing-]{dualizing}
\externaldocument[duality-]{duality}
\externaldocument[discriminant-]{discriminant}
\externaldocument[derham-]{derham}
\externaldocument[local-cohomology-]{local-cohomology}
\externaldocument[algebraization-]{algebraization}
\externaldocument[curves-]{curves}
\externaldocument[resolve-]{resolve}
\externaldocument[models-]{models}
\externaldocument[functors-]{functors}
\externaldocument[equiv-]{equiv}
\externaldocument[pione-]{pione}
\externaldocument[etale-cohomology-]{etale-cohomology}
\externaldocument[proetale-]{proetale}
\externaldocument[relative-cycles-]{relative-cycles}
\externaldocument[more-etale-]{more-etale}
\externaldocument[trace-]{trace}
\externaldocument[crystalline-]{crystalline}
\externaldocument[spaces-]{spaces}
\externaldocument[spaces-properties-]{spaces-properties}
\externaldocument[spaces-morphisms-]{spaces-morphisms}
\externaldocument[decent-spaces-]{decent-spaces}
\externaldocument[spaces-cohomology-]{spaces-cohomology}
\externaldocument[spaces-limits-]{spaces-limits}
\externaldocument[spaces-divisors-]{spaces-divisors}
\externaldocument[spaces-over-fields-]{spaces-over-fields}
\externaldocument[spaces-topologies-]{spaces-topologies}
\externaldocument[spaces-descent-]{spaces-descent}
\externaldocument[spaces-perfect-]{spaces-perfect}
\externaldocument[spaces-more-morphisms-]{spaces-more-morphisms}
\externaldocument[spaces-flat-]{spaces-flat}
\externaldocument[spaces-groupoids-]{spaces-groupoids}
\externaldocument[spaces-more-groupoids-]{spaces-more-groupoids}
\externaldocument[bootstrap-]{bootstrap}
\externaldocument[spaces-pushouts-]{spaces-pushouts}
\externaldocument[spaces-chow-]{spaces-chow}
\externaldocument[groupoids-quotients-]{groupoids-quotients}
\externaldocument[spaces-more-cohomology-]{spaces-more-cohomology}
\externaldocument[spaces-simplicial-]{spaces-simplicial}
\externaldocument[spaces-duality-]{spaces-duality}
\externaldocument[formal-spaces-]{formal-spaces}
\externaldocument[restricted-]{restricted}
\externaldocument[spaces-resolve-]{spaces-resolve}
\externaldocument[formal-defos-]{formal-defos}
\externaldocument[defos-]{defos}
\externaldocument[cotangent-]{cotangent}
\externaldocument[examples-defos-]{examples-defos}
\externaldocument[algebraic-]{algebraic}
\externaldocument[examples-stacks-]{examples-stacks}
\externaldocument[stacks-sheaves-]{stacks-sheaves}
\externaldocument[criteria-]{criteria}
\externaldocument[artin-]{artin}
\externaldocument[quot-]{quot}
\externaldocument[stacks-properties-]{stacks-properties}
\externaldocument[stacks-morphisms-]{stacks-morphisms}
\externaldocument[stacks-limits-]{stacks-limits}
\externaldocument[stacks-cohomology-]{stacks-cohomology}
\externaldocument[stacks-perfect-]{stacks-perfect}
\externaldocument[stacks-introduction-]{stacks-introduction}
\externaldocument[stacks-more-morphisms-]{stacks-more-morphisms}
\externaldocument[stacks-geometry-]{stacks-geometry}
\externaldocument[moduli-]{moduli}
\externaldocument[moduli-curves-]{moduli-curves}
\externaldocument[examples-]{examples}
\externaldocument[exercises-]{exercises}
\externaldocument[guide-]{guide}
\externaldocument[desirables-]{desirables}
\externaldocument[coding-]{coding}
\externaldocument[obsolete-]{obsolete}
\externaldocument[fdl-]{fdl}
\externaldocument[index-]{index}

% Theorem environments.
%
\theoremstyle{plain}
\newtheorem{theorem}[subsection]{Theorem}
\newtheorem{proposition}[subsection]{Proposition}
\newtheorem{lemma}[subsection]{Lemma}

\theoremstyle{definition}
\newtheorem{definition}[subsection]{Definition}
\newtheorem{example}[subsection]{Example}
\newtheorem{exercise}[subsection]{Exercise}
\newtheorem{situation}[subsection]{Situation}

\theoremstyle{remark}
\newtheorem{remark}[subsection]{Remark}
\newtheorem{remarks}[subsection]{Remarks}

\numberwithin{equation}{subsection}

% Macros
%
\def\lim{\mathop{\mathrm{lim}}\nolimits}
\def\colim{\mathop{\mathrm{colim}}\nolimits}
\def\Spec{\mathop{\mathrm{Spec}}}
\def\Hom{\mathop{\mathrm{Hom}}\nolimits}
\def\Ext{\mathop{\mathrm{Ext}}\nolimits}
\def\SheafHom{\mathop{\mathcal{H}\!\mathit{om}}\nolimits}
\def\SheafExt{\mathop{\mathcal{E}\!\mathit{xt}}\nolimits}
\def\Sch{\mathit{Sch}}
\def\Mor{\mathop{\mathrm{Mor}}\nolimits}
\def\Ob{\mathop{\mathrm{Ob}}\nolimits}
\def\Sh{\mathop{\mathit{Sh}}\nolimits}
\def\NL{\mathop{N\!L}\nolimits}
\def\CH{\mathop{\mathrm{CH}}\nolimits}
\def\proetale{{pro\text{-}\acute{e}tale}}
\def\etale{{\acute{e}tale}}
\def\QCoh{\mathit{QCoh}}
\def\Ker{\mathop{\mathrm{Ker}}}
\def\Im{\mathop{\mathrm{Im}}}
\def\Coker{\mathop{\mathrm{Coker}}}
\def\Coim{\mathop{\mathrm{Coim}}}

% Boxtimes
%
\DeclareMathSymbol{\boxtimes}{\mathbin}{AMSa}{"02}

%
% Macros for moduli stacks/spaces
%
\def\QCohstack{\mathcal{QC}\!\mathit{oh}}
\def\Cohstack{\mathcal{C}\!\mathit{oh}}
\def\Spacesstack{\mathcal{S}\!\mathit{paces}}
\def\Quotfunctor{\mathrm{Quot}}
\def\Hilbfunctor{\mathrm{Hilb}}
\def\Curvesstack{\mathcal{C}\!\mathit{urves}}
\def\Polarizedstack{\mathcal{P}\!\mathit{olarized}}
\def\Complexesstack{\mathcal{C}\!\mathit{omplexes}}
% \Pic is the operator that assigns to X its picard group, usage \Pic(X)
% \Picardstack_{X/B} denotes the Picard stack of X over B
% \Picardfunctor_{X/B} denotes the Picard functor of X over B
\def\Pic{\mathop{\mathrm{Pic}}\nolimits}
\def\Picardstack{\mathcal{P}\!\mathit{ic}}
\def\Picardfunctor{\mathrm{Pic}}
\def\Deformationcategory{\mathcal{D}\!\mathit{ef}}
\usepackage{cleveref}
\crefname{originalreference}{original reference}{original references}
\Crefname{originalreference}{Original reference}{Original references}
\makeatletter
\newlabel{algebra-lemma-glue-modules}{{00EQ}{source}{}{sourcepointer-1}{}}
\newlabel{algebra-lemma-glue-modules@cref}{{[originalreference][0][]00EQ}{[1][0][]source}{}{}{}}
\makeatother
\crefname{originallemma}{original lemma}{original lemmas}
\Crefname{originallemma}{Original lemma}{Original lemmas}
\makeatletter
\newlabel{lemma-compare-constructions}{{01I7}{source}{}{sourcepointer-2}{}}
\newlabel{lemma-compare-constructions@cref}{{[originallemma][0][]01I7}{[1][0][]source}{}{}{}}
\makeatother
\crefname{originallemma}{original lemma}{original lemmas}
\Crefname{originallemma}{Original lemma}{Original lemmas}
\makeatletter
\newlabel{lemma-standard-open}{{01HS}{source}{}{sourcepointer-3}{}}
\newlabel{lemma-standard-open@cref}{{[originallemma][0][]01HS}{[1][0][]source}{}{}{}}
\makeatother
\crefname{originalreference}{original reference}{original references}
\Crefname{originalreference}{Original reference}{Original references}
\makeatletter
\newlabel{modules-lemma-quasi-coherent-module}{{01BK}{source}{}{sourcepointer-4}{}}
\newlabel{modules-lemma-quasi-coherent-module@cref}{{[originalreference][0][]01BK}{[1][0][]source}{}{}{}}
\makeatother
\crefname{originalreference}{original reference}{original references}
\Crefname{originalreference}{Original reference}{Original references}
\makeatletter
\newlabel{sheaves-lemma-mapping-property-glue}{{00AN}{source}{}{sourcepointer-5}{}}
\newlabel{sheaves-lemma-mapping-property-glue@cref}{{[originalreference][0][]00AN}{[1][0][]source}{}{}{}}
\makeatother
\crefname{originalreference}{original reference}{original references}
\Crefname{originalreference}{Original reference}{Original references}
\makeatletter
\newlabel{sheaves-section-glueing-sheaves}{{00AK}{source}{}{sourcepointer-6}{}}
\newlabel{sheaves-section-glueing-sheaves@cref}{{[originalreference][0][]00AK}{[1][0][]source}{}{}{}}
\makeatother
\crefname{originalreference}{original reference}{original references}
\Crefname{originalreference}{Original reference}{Original references}
\makeatletter
\newlabel{topology-definition-locally-quasi-compact}{{0068}{source}{}{sourcepointer-7}{}}
\newlabel{topology-definition-locally-quasi-compact@cref}{{[originalreference][0][]0068}{[1][0][]source}{}{}{}}
\makeatother
\title{Quasi-coherent sheaves on an affine scheme come from their global sections}
\author{The Stacks Project Authors}
\date{stacks-project:schemes.tex:local-snapshot-20261012; GNU FDL 1.2 or later}
\begin{document}
\maketitle
\noindent\textit{Editorial note.} Counted span is the author's own proof environment, source lines 1504-1604, following the statement at 1495-1502.\par\medskip
\noindent\textit{Editorial note.} The sentence 'We leave it to the reader to show that m -> s is a R-module map' (\textasciitilde{}1595-1596) is a mechanical linearity check after the map is displayed explicitly; flagged in the notes.\par\medskip
\noindent\textit{Editorial note.} The triple-overlap cocycle 'It is clear' (\textasciitilde{}1534-1535) is read off the constructive definition of psi\_ij as a composition of restrictions of isomorphisms.\par\medskip
\noindent\textit{Allowed standard prerequisites (editorial).}\begin{itemize}
\item Local representation of quasi-coherent modules. Modules, Lemma modules-lemma-quasi-coherent-module (cross-chapter, cited at \textasciitilde{}1510). Gives a qc-neighborhood cover on which F is a tilde sheaf.
\item Standard open refinement. lemma-standard-open (same chapter, cited at \textasciitilde{}1516). Refines the neighborhood cover to a standard open cover D(f\_i).
\item Comparison of the two tilde constructions. lemma-compare-constructions (same chapter, cited at \textasciitilde{}1548, \textasciitilde{}1551, \textasciitilde{}1597). Identifies tilde sheaves with their modules on affines and converts sheaf maps to module maps.
\item Glueing of modules. Algebra, Lemma algebra-lemma-glue-modules (cross-chapter, cited at \textasciitilde{}1572). Produces M with M\_i = M\_\{f\_i\} compatible with the psi\_ij.
\item Glueing sheaves from data. Sheaves, Section sheaves-section-glueing-sheaves, in particular sheaves-lemma-mapping-property-glue (cross-chapter, cited at \textasciitilde{}1580-1582). Frames F and tilde(M) as glueing data over the covering.
\end{itemize}\medskip
\ifdefined\hypertarget\hypertarget{sourcepointer-1}{}\fi
\noindent\textit{Source reference.} Original reference 00EQ: focused-seeds/geometry/S9-schemes-quasi-coherent-affine/ref-tags.txt, source lines 5--5. Navigation pointer only: the counted proof cites this statement by \textbackslash{}ref as a standing input; the deduction the counted proof has to make is carried out inside the counted span, and no part of that deduction is delegated to this pointer.\par\medskip
\ifdefined\hypertarget\hypertarget{sourcepointer-2}{}\fi
\noindent\textit{Source reference.} Original lemma 01I7: sources/stacks-schemes/pinned-original-source.tex, source lines 1295--1295. Navigation pointer only: the counted proof cites this statement by \textbackslash{}ref as a standing input; the deduction the counted proof has to make is carried out inside the counted span, and no part of that deduction is delegated to this pointer.\par\medskip
\ifdefined\hypertarget\hypertarget{sourcepointer-3}{}\fi
\noindent\textit{Source reference.} Original lemma 01HS: sources/stacks-schemes/pinned-original-source.tex, source lines 737--737. Navigation pointer only: the counted proof cites this statement by \textbackslash{}ref as a standing input; the deduction the counted proof has to make is carried out inside the counted span, and no part of that deduction is delegated to this pointer.\par\medskip
\ifdefined\hypertarget\hypertarget{sourcepointer-4}{}\fi
\noindent\textit{Source reference.} Original reference 01BK: focused-seeds/geometry/S9-schemes-quasi-coherent-affine/ref-tags.txt, source lines 6--6. Navigation pointer only: the counted proof cites this statement by \textbackslash{}ref as a standing input; the deduction the counted proof has to make is carried out inside the counted span, and no part of that deduction is delegated to this pointer.\par\medskip
\ifdefined\hypertarget\hypertarget{sourcepointer-5}{}\fi
\noindent\textit{Source reference.} Original reference 00AN: focused-seeds/geometry/S9-schemes-quasi-coherent-affine/ref-tags.txt, source lines 7--7. Navigation pointer only: the counted proof cites this statement by \textbackslash{}ref as a standing input; the deduction the counted proof has to make is carried out inside the counted span, and no part of that deduction is delegated to this pointer.\par\medskip
\ifdefined\hypertarget\hypertarget{sourcepointer-6}{}\fi
\noindent\textit{Source reference.} Original reference 00AK: focused-seeds/geometry/S9-schemes-quasi-coherent-affine/ref-tags.txt, source lines 8--8. Navigation pointer only: the counted proof cites this statement by \textbackslash{}ref as a standing input; the deduction the counted proof has to make is carried out inside the counted span, and no part of that deduction is delegated to this pointer.\par\medskip
\ifdefined\hypertarget\hypertarget{sourcepointer-7}{}\fi
\noindent\textit{Source reference.} Original reference 0068: focused-seeds/geometry/S9-schemes-quasi-coherent-affine/ref-tags.txt, source lines 9--9. Navigation pointer only: the counted proof cites this statement by \textbackslash{}ref as a standing input; the deduction the counted proof has to make is carried out inside the counted span, and no part of that deduction is delegated to this pointer.\par\medskip
\begin{lemma}
\label{lemma-quasi-coherent-affine}
Let $(X, \mathcal{O}_X) = (\Spec(R), \mathcal{O}_{\Spec(R)})$
be an affine scheme. Let $\mathcal{F}$ be a
quasi-coherent $\mathcal{O}_X$-module. Then
$\mathcal{F}$ is isomorphic to the sheaf associated to
the $R$-module $\Gamma(X, \mathcal{F})$.
\end{lemma}

\begin{proof}
Let $\mathcal{F}$ be a quasi-coherent $\mathcal{O}_X$-module.
Since every standard open $D(f)$ is quasi-compact we see that
$X$ is locally quasi-compact, i.e., every point has a fundamental
system of quasi-compact neighbourhoods, see Topology,
Definition \ref{topology-definition-locally-quasi-compact}.
Hence by Modules, Lemma \ref{modules-lemma-quasi-coherent-module}
for every prime $\mathfrak p \subset R$ corresponding to $x \in X$
there exists an open neighbourhood $x \in U \subset X$ such that
$\mathcal{F}|_U$ is isomorphic to the quasi-coherent
sheaf associated to some $\mathcal{O}_X(U)$-module $M$.
In other words, we get an open covering by $U$'s with this property.
By Lemma \ref{lemma-standard-open} for example we can refine this
covering to  a standard open covering.
Thus we get a covering $\Spec(R) = \bigcup D(f_i)$
and $R_{f_i}$-modules $M_i$ and isomorphisms
$\varphi_i : \mathcal{F}|_{D(f_i)} \to \mathcal{F}_{M_i}$
for some $R_{f_i}$-module $M_i$. On the overlaps
we get isomorphisms
$$
\xymatrix{
\mathcal{F}_{M_i}|_{D(f_if_j)}
\ar[rr]^{\varphi_i^{-1}|_{D(f_if_j)}}
& &
\mathcal{F}|_{D(f_if_j)}
\ar[rr]^{\varphi_j|_{D(f_if_j)}}
& &
\mathcal{F}_{M_j}|_{D(f_if_j)}.
}
$$
Let us denote these $\psi_{ij}$. It is clear that
we have the cocycle condition
$$
\psi_{jk}|_{D(f_if_jf_k)}
\circ
\psi_{ij}|_{D(f_if_jf_k)}
=
\psi_{ik}|_{D(f_if_jf_k)}
$$
on triple overlaps.

\medskip\noindent
Recall that each of the open subspaces $D(f_i)$, $D(f_if_j)$,
$D(f_if_jf_k)$ is an affine scheme. Hence the sheaves $\mathcal{F}_{M_i}$
are isomorphic to the sheaves $\widetilde M_i$ by Lemma
\ref{lemma-compare-constructions} above. In particular we see that
$\mathcal{F}_{M_i}(D(f_if_j)) = (M_i)_{f_j}$, etc.
Also by Lemma \ref{lemma-compare-constructions} above we see
that $\psi_{ij}$ corresponds to a unique $R_{f_if_j}$-module isomorphism
$$
\psi_{ij} : (M_i)_{f_j} \longrightarrow (M_j)_{f_i}
$$
namely, the effect of $\psi_{ij}$ on sections over $D(f_if_j)$.
Moreover these then satisfy the cocycle condition that
$$
\xymatrix{
(M_i)_{f_jf_k}
\ar[rd]_{\psi_{ij}}
\ar[rr]^{\psi_{ik}}
& &
(M_k)_{f_if_j} \\
&
(M_j)_{f_if_k} \ar[ru]_{\psi_{jk}}
}
$$
commutes (for any triple $i, j, k$).

\medskip\noindent
Now Algebra, Lemma \ref{algebra-lemma-glue-modules}
shows that there exist an $R$-module $M$ such that
$M_i = M_{f_i}$ compatible with the morphisms $\psi_{ij}$.
Consider $\mathcal{F}_M = \widetilde M$. At this point it is
a formality to show that $\widetilde M$ is isomorphic to
the quasi-coherent sheaf $\mathcal{F}$ we started out with.
Namely, the sheaves $\mathcal{F}$ and $\widetilde M$ give
rise to isomorphic sets of glueing data of sheaves of $\mathcal{O}_X$-modules
with respect to the covering $X = \bigcup D(f_i)$, see
Sheaves, Section \ref{sheaves-section-glueing-sheaves}
and in particular Lemma \ref{sheaves-lemma-mapping-property-glue}.
Explicitly, in the current situation, this boils down to
the following argument: Let us construct an $R$-module map
$$
M \longrightarrow \Gamma(X, \mathcal{F}).
$$
Namely, given $m \in M$ we get $m_i = m/1 \in M_{f_i} = M_i$
by construction of $M$. By construction of $M_i$ this corresponds
to a section $s_i \in \mathcal{F}(U_i)$. (Namely, $\varphi^{-1}_i(m_i)$.)
We claim that $s_i|_{D(f_if_j)} = s_j|_{D(f_if_j)}$. This is
true because, by construction of $M$, we have $\psi_{ij}(m_i) = m_j$,
and by the construction of the $\psi_{ij}$. By the sheaf condition of
$\mathcal{F}$ this collection of sections gives rise to a unique
section $s$ of $\mathcal{F}$ over $X$. We leave it to the reader
to show that $m \mapsto s$ is a $R$-module map.
By Lemma \ref{lemma-compare-constructions} we obtain an associated
$\mathcal{O}_X$-module map
$$
\widetilde M \longrightarrow \mathcal{F}.
$$
By construction this map reduces to the isomorphisms
$\varphi_i^{-1}$ on each $D(f_i)$ and hence is an isomorphism.
\end{proof}

\end{document}
